%! Sinusoidal Function Transformations — Unit 6, Lesson 6.7 — Guided Notes (Answer Key)
\documentclass[10pt]{article}
\usepackage{precalculus-article}
\usepackage{precalculus-key}   % pulls in -boxes; do NOT also load -boxes
\newcommand{\TallMath}[1]{$\displaystyle #1\rule[-1.4em]{0pt}{3.2em}$}

\begin{document}

\pageheader{Unit 6, Lesson 6.7}{Guided Notes --- Answer Key}
\namedateperiod

\vspace{0.10in}

%----------------------------------------------------------------------
% Objectives
%----------------------------------------------------------------------
\begin{objectivebox}
\textbf{I will be able to\ldots}
\begin{enumerate}[leftmargin=1.5em, itemsep=4pt]
  \item \textbf{LO 3.6.A} --- Identify the amplitude, vertical shift (midline),
    period, and phase shift of a sinusoidal function from its equation.
    \ans{Identify $|a|$, $d$, $2\pi/|b|$, and $-c$ from $a\sin(b(\theta+c))+d$.}
  \item \textbf{LO 3.6.A} --- Explain how each parameter in $f(\theta)=a\sin(b(\theta+c))+d$
    corresponds to an additive or multiplicative transformation of $y=\sin\theta$.
    \ans{$a$ dilates vertically; $d$ shifts vertically; $b$ dilates horizontally; $c$ shifts horizontally.}
\end{enumerate}
\end{objectivebox}

\vspace{0.08in}

%----------------------------------------------------------------------
% Vocabulary
%----------------------------------------------------------------------
\begin{vocabbox}
\begin{description}[leftmargin=0pt, itemsep=6pt]
  \item[\termblanklong{sinusoidal function}]
    A function equivalent to $a\sin(b(\theta+c))+d$ (or the cosine form) with $a \neq 0$.
    A \ans{transformation} of the parent sine or cosine function.

  \item[\termblanklong{amplitude}]
    Equals $|a|$ --- half the vertical distance from the
    \ans{midline} to the \ans{maximum (or minimum)}.
    Controls the vertical \ans{stretch or compression}.

  \item[\termblanklong{midline}]
    The horizontal line $y = $ \ans{$d$}. The vertical
    \ans{shift} of the function. Average of max and min outputs.

  \item[\termblanklong{period}]
    Length of one complete cycle; $T = \dfrac{2\pi}{\ans{|b|}}$.
    Controlled by the horizontal dilation factor $b$.

  \item[\termblanklong{phase shift}]
    Horizontal translation equal to \ans{$-c$} when the function
    is written in factored form $b(\theta + c)$.
    Positive $c$ shifts \ans{left}; negative $c$ shifts
    \ans{right}.
\end{description}
\end{vocabbox}

\vspace{0.08in}

%----------------------------------------------------------------------
% Hook
%----------------------------------------------------------------------
\begin{hookbox}
\textbf{Entry Question:} (Graph as in blank notes.)

\vspace{4pt}
\textbf{Q1.} Graph A has amplitude \ans{1} and midline $y = $ \ans{0}.\\[4pt]
\textbf{Q2.} Graph B has amplitude \ans{3} and midline $y = $ \ans{1}.\\[4pt]
\textbf{Q3.} What equation could describe Graph B? $y = $ \ans{$3\sin\theta + 1$}

\end{hookbox}

\vspace{0.08in}

%----------------------------------------------------------------------
% LO 3.6.A (Part 1) — Vertical Transformations
%----------------------------------------------------------------------
\begin{notesbox}{LO 3.6.A --- Part 1: Vertical Transformations ($a$ and $d$) --- Key}

\textbf{Vertical Dilation:} $g(\theta) = a\sin(\theta)$

Multiplying all outputs by $a$ changes the \ans{amplitude}. The new amplitude is \ans{$|a|$}.
If $a < 0$, the graph is also \ans{reflected (flipped)} over the midline.

\vspace{6pt}
\renewcommand{\arraystretch}{1.5}
\begin{tabular}{|l|c|c|c|c|}
\hline
\textbf{Function} & \textbf{Amplitude} & \textbf{Max} & \textbf{Min} & \textbf{Midline} \\
\hline
$y = \sin\theta$ & \ans{1} & \ans{1} & \ans{$-1$} & \ans{$y=0$} \\
\hline
$y = 4\sin\theta$ & \ans{4} & \ans{4} & \ans{$-4$} & \ans{$y=0$} \\
\hline
$y = -2\sin\theta$ & \ans{2} & \ans{2} & \ans{$-2$} & \ans{$y=0$} \\
\hline
$y = \tfrac{1}{2}\sin\theta$ & \ans{$\tfrac{1}{2}$} & \ans{$\tfrac{1}{2}$} & \ans{$-\tfrac{1}{2}$} & \ans{$y=0$} \\
\hline
\end{tabular}

\vspace{8pt}
\textbf{Vertical Translation:} $g(\theta) = \sin(\theta) + d$

Adding $d$ shifts the \ans{midline} to $y = d$.
The amplitude is \ans{not affected / unchanged}.

\vspace{6pt}
\begin{tabular}{|l|c|c|c|c|}
\hline
\textbf{Function} & \textbf{Amplitude} & \textbf{Max} & \textbf{Min} & \textbf{Midline} \\
\hline
$y = \sin\theta + 3$ & \ans{1} & \ans{4} & \ans{2} & \ans{$y=3$} \\
\hline
$y = 2\sin\theta - 1$ & \ans{2} & \ans{1} & \ans{$-3$} & \ans{$y=-1$} \\
\hline
\end{tabular}

\end{notesbox}

\vspace{0.08in}

%----------------------------------------------------------------------
% LO 3.6.A (Part 2) — Horizontal Transformations
%----------------------------------------------------------------------
\begin{notesbox}{LO 3.6.A --- Part 2: Horizontal Transformations ($b$ and $c$) --- Key}

\textbf{Horizontal Dilation:} $g(\theta) = \sin(b\theta)$

\vspace{4pt}
If $|b| > 1$: the graph is \ans{compressed horizontally} (period is shorter). \\
If $|b| < 1$: the graph is \ans{stretched horizontally} (period is longer).

\vspace{6pt}
\begin{tabular}{|l|c|c|}
\hline
\textbf{Function} & \textbf{$b$ value} & \textbf{Period} \\
\hline
$y = \sin(3\theta)$ & \ans{3} & \ans{$2\pi/3$} \\
\hline
$y = \sin\!\left(\tfrac{1}{2}\theta\right)$ & \ans{$1/2$} & \ans{$4\pi$} \\
\hline
$y = \sin(4\theta)$ & \ans{4} & \ans{$\pi/2$} \\
\hline
\end{tabular}

\vspace{8pt}
Phase shift $= $ \ans{$-c$} (use the factored form to read it correctly).

\vspace{4pt}
\textbf{Important --- Factoring first!}
$$\sin(2\theta + \pi) = \sin\!\bigl(2(\theta + \ans{$\pi/2$})\bigr)$$
Phase shift $= -\ans{$\pi/2$} = $ \ans{$-\pi/2$ (left $\pi/2$ units)}

\end{notesbox}

\vspace{0.08in}

%----------------------------------------------------------------------
% LO 3.6.A (Part 3) — General Form and Summary
%----------------------------------------------------------------------
\begin{notesbox}{LO 3.6.A --- Part 3: The General Form --- Key}

\textbf{General Sinusoidal Function:}
$$f(\theta) = a\sin\bigl(b(\theta + c)\bigr) + d$$

\renewcommand{\arraystretch}{1.6}
\begin{tabularx}{\linewidth}{@{} p{0.14\linewidth} p{0.26\linewidth} X @{}}
\hline
\textbf{Parameter} & \textbf{Feature} & \textbf{Formula} \\
\hline
$a$ & Amplitude & $|a|$ \quad (sign determines reflection) \\
$b$ & Period & $T = \dfrac{2\pi}{|b|}$ \\
$c$ & Phase shift & $-c$ \quad (left if $c > 0$, right if $c < 0$) \\
$d$ & Midline & $y = d$ \\
\hline
\end{tabularx}

\vspace{8pt}
\textbf{Worked Example:} $f(\theta) = -2\sin\!\bigl(\tfrac{1}{2}(\theta - \pi)\bigr) + 3$

\vspace{4pt}
\begin{tabular}{ll}
Amplitude: & \ans{$|-2| = 2$} \\[6pt]
Period: & \ans{$2\pi / (1/2) = 4\pi$} \\[6pt]
Phase shift: & \ans{$-(-\pi) = \pi$, shifted right $\pi$ units} \\[6pt]
Midline: & \ans{$y = 3$} \\[6pt]
Reflected? & \ans{Yes, because $a = -2 < 0$} \\
\end{tabular}

\vspace{8pt}
\textbf{Sine--Cosine Equivalence:}
$$2\cos\theta = 2\sin\!\left(\theta + \ans{$\dfrac{\pi}{2}$}\right)$$

\end{notesbox}

\vspace{0.08in}

%----------------------------------------------------------------------
% Practice
%----------------------------------------------------------------------
\begin{practicebox}

\textbf{Guided Practice 1.}

\vspace{4pt}
\renewcommand{\arraystretch}{1.8}
\begin{tabular}{|l|c|c|c|c|}
\hline
\textbf{Function} & \textbf{Amplitude} & \textbf{Period} & \textbf{Phase Shift} & \textbf{Midline} \\
\hline
$f(\theta) = 5\sin(3\theta) - 2$ & \ans{5} & \ans{$2\pi/3$} & \ans{none (0)} & \ans{$y=-2$} \\
\hline
$g(\theta) = -\cos\!\left(\tfrac{\pi}{2}\theta\right) + 4$ & \ans{1} & \ans{4} & \ans{none (0)} & \ans{$y=4$} \\
\hline
$h(\theta) = 3\sin\!\left(2\left(\theta + \dfrac{\pi}{6}\right)\right)$ & \ans{3} & \ans{$\pi$} & \ans{$-\pi/6$ (left)} & \ans{$y=0$} \\
\hline
\end{tabular}

\vspace{0.10in}
\textbf{Guided Practice 2.}

\vspace{4pt}
(a)\quad $3\theta - \pi = 3\!\left(\theta - \ans{$\pi/3$}\right)$

\vspace{4pt}
(b)\quad Phase shift $= $ \ans{$\pi/3$ to the right}

\vspace{4pt}
(c)\quad Amplitude \ans{4}, period \ans{$2\pi/3$}, phase shift \ans{$\pi/3$ right}, midline \ans{$y=1$}.

\end{practicebox}

\end{document}
